\subsection{The canonical neighborhood obstruction for algebraic-dimension-zero Hopf components}

\begin{statement}{C.2.UN.38}{Proved}{Every curve of a smooth minimal algebraic-dimension-zero Hopf normalization lies in the fibre boundary}
Let $f:U\to\Delta$ be a proper degeneration of
connected complex two-tori from a smooth threefold,
and write its full central fibre as $\sum_i m_iS_i$.
On the open subset of $S_i$ where the reduced
support is smooth, its canonical line has an
actual trivial $m_i$-th power in a neighborhood
of every compact subset. In particular, a
smooth compact elliptic curve $C$ in that open
subset satisfies
\[
                         N_{C/S_i}^{\otimes m_i}
                                      \simeq\OO_C.    \tag{C.HB.1}
\]

Suppose the normalization of a central component
is itself a smooth minimal Hopf surface of
algebraic dimension zero. Then every elliptic
curve on that normalization is contained in
the inverse image of the singular locus of
the reduced support. For a diagonal primary
Hopf surface this includes both axes. For a
non-diagonal primary Hopf surface it includes
its unique axis; in this case the conclusion
uses the full canonical neighborhood and
the explicit higher-jet calculation below.
The same conclusion descends through every
actual finite primary Hopf covering.
\end{statement}

\begin{proof}
The proper direct image $f_*K_U$ is a
torsion-free coherent sheaf of generic rank
one on the disc and hence is a line bundle.
Choose a generator after shrinking the disc.
Its evaluation is a holomorphic section
$\omega$ of $K_U$. On each nonzero fibre,
proper base change identifies this evaluation
with a nonzero canonical form on that torus,
up to the fixed base differential. It has
no zero there. Consequently its exact divisor
is vertical:
\[
                   \operatorname{div}(\omega)
                               =\sum_i k_iS_i,
                        \qquad k_i\ge0.              \tag{C.HB.2}
\]
Let $W\subset U$ be a neighborhood of a
compact subset in the indicated regular
part of $S=S_i$, small enough to avoid all
other components and the singularities of
$S$. On $W$ the actual fibre equation has
divisor $mS$, where $m=m_i$. Thus it gives
the trivialization of $\OO_W(mS)$, and
(C.HB.2) gives
\[
 \begin{gathered}
 K_U|_W\simeq\OO_W(k_iS),\\
 K_S|_{S\cap W}\simeq
                  \OO_W((k_i+1)S)|_{S\cap W},\\
 (K_S|_{S\cap W})^{\otimes m}\simeq\OO_{S\cap W}.
 \end{gathered}
                                                               \tag{C.HB.3}
\]
These are line-bundle isomorphisms on the
neighborhood, obtained from the actual
canonical and fibre sections. Adjunction
for an elliptic $C\subset S\cap W$ gives
$N_{C/S}=K_C\otimes K_S^{-1}|_C$.
Since $K_C=\OO_C$, (C.HB.1) follows.

Consider first the diagonal primary
Hopf surface
\[
 H_{\alpha,\beta}=(\C^2\setminus\{0\})/
       \langle(z,w)\mapsto(\alpha z,\beta w)\rangle,
                    \qquad0<|\alpha|,|\beta|<1.
                                                               \tag{C.HB.4}
\]
The axis $C_\alpha$ is $\C^*/\alpha^{\Z}$,
and its actual normal line is the quotient
of $\C^*\times\C$ by
$(z,v)\mapsto(\alpha z,\beta v)$.
If a positive tensor power $n$ of this
line is trivial, a nowhere-zero holomorphic
function $a$ on $\C^*$ satisfies
$a(\alpha z)=\beta^n a(z)$.
In its Laurent expansion each nonzero
coefficient of $z^j$ requires
$\alpha^j=\beta^n$. Since $|\alpha|<1$,
there is at most one such exponent, and
there must be one. Its exponent is positive
because $|\beta|<1$ and $n>0$.
Conversely that equality gives the
explicit trivialization $a(z)=z^j$.
It also gives the nonconstant invariant
meromorphic function $w^n/z^j$ on the
Hopf surface. Algebraic dimension zero
therefore makes the normal of each axis
non-torsion. Formula (C.HB.1) excludes
an entire such axis from the regular
part of the reduced support.

The non-diagonal primary normal form
requires a further argument. The primary
normal forms and finite unramified cover
used here are the explicit ones recalled
in \cite[Theorem~2.1]{Khaled2024Hopf}
and \cite[\S1.1]{Nakamura2008VII}.
Write the non-diagonal contraction
as
\[
 H=(\C^2\setminus\{0\})/\langle F\rangle,
 \qquad F(z,w)=(\beta^r z+w^r,\beta w),
                  \quad r\ge1,\quad0<|\beta|<1,
                                                               \tag{C.HB.5}
\]
and let $C$ be the image of $w=0$.
The invariant meromorphic two-form
$dz\wedge dw/w^{r+1}$ gives the actual
canonical divisor
\[
                         K_H\simeq\OO_H(-(r+1)C).
                                                               \tag{C.HB.6}
\]
Its normal line has torsion, so (C.HB.1)
alone gives no contradiction. Instead,
(C.HB.3) on a neighborhood of $C$ would
make $\OO(m(r+1)C)$ trivial there.
Its canonical divisor section would
then be a holomorphic function with
positive divisor $m(r+1)C$ on that
neighborhood. We show directly that
no such function exists.

Lift it to the invariant neighborhood
of $\C^*\times\{0\}$ in the covering
space and expand in the normal variable:
\[
 h(z,w)=\sum_{l\ge j}h_l(z)w^l,
       \qquad j>0,\qquad h_j\ne0,
       \qquad h_l\in\OO(\C^*),
       \qquad h\circ F=h.                           \tag{C.HB.7}
\]
The radius of convergence may depend
on $z$; the Taylor coefficients are
nevertheless holomorphic on all of
$\C^*$ because the lifted neighborhood
contains that entire axis. The leading
coefficient satisfies
$\beta^j h_j(\beta^r z)=h_j(z)$.
Its Laurent expansion forces
\[
                      j=rk,\qquad
                      h_j(z)=c z^{-k},
                      \qquad k\ge1,\quad c\ne0.
                                                               \tag{C.HB.8}
\]
At order $j+r$, the same invariant
function equation is exactly
\[
 \beta^{j+r}h_{j+r}(\beta^r z)-h_{j+r}(z)
                    =k c\beta^{-r}z^{-k-1}.          \tag{C.HB.9}
\]
Indeed the only extra contribution
at that order is the first derivative
of $h_j(\beta^r z+w^r)\beta^j$;
all other Taylor terms start at
higher order. In the Laurent
expansion of the left side, the
coefficient multiplying $z^{-k-1}$
is zero, since
$\beta^{j+r+r(-k-1)}-1=0$.
The right side has the nonzero
coefficient $k c\beta^{-r}$.
This contradiction proves the
neighborhood obstruction for
the non-diagonal axis as well.

Finally let $\nu:H\to S$ be the
given smooth minimal Hopf
normalization. Its conductor
and the inverse images of
the other central components
have curve support. Outside
this support $\nu$ is the
actual isomorphism onto the
regular locus of the reduced
fibre. For clarity, there
is no extra isolated
normalization defect: a
hypersurface is Cohen--Macaulay,
so its nonnormal locus is
the conductor curve locus;
where it is normal, the
normalization is an isomorphism,
and smoothness of $H$ then
gives smoothness there.
An intersection with another
central component is a
Cartier curve divisor and
its finite pullback also
has curve support.

We also verify the curve assertion
used in applying the obstruction.
The inverse image of a compact
curve in the covering punctured
plane extends analytically across
the origin by the removable-set
theorem for analytic sets, since
its dimension is one and the
added set has dimension zero.
Its germ has a reduced equation
$a$ satisfying $a\circ F=v a$
for a holomorphic unit $v$.
After shrinking to a contracting
ball, the product solving the
unit cocycle converges: if
$u=v/v(0)$, then
$b=\prod_{j\ge0}u\circ F^j$
is a holomorphic unit and
$b\circ F=b/u$.
Thus $h=ba$ satisfies
$h\circ F=v(0)h$.
For the diagonal contraction,
two distinct nonzero monomials
in $h$ would give
$\alpha^i\beta^j=1$ for a
nonzero integer pair; because
both moduli are less than one,
this is a positive-power
relation between $\alpha$
and $\beta$ and gives a
nonconstant meromorphic
invariant. In algebraic
dimension zero there is
therefore just one monomial,
so the curve consists of
axes. For (C.HB.5), decompose
the convergent power series
into weighted degrees with
weights $r,1$ on $z,w$.
The contraction preserves
each finite-dimensional
weighted piece and acts
there with sole eigenvalue
$\beta^d$ in degree $d$.
Their distinct moduli force
an eigenfunction into one
degree. Write it as
$w^dP(z/w^r)$; after taking
out $\beta^d$, invariance
requires the polynomial
$P$ to be periodic under
translation by $\beta^{-r}$.
It is constant. Thus only
the axis $w=0$ occurs.
Analytic continuation along
the contraction orbits
gives this assertion on
the entire punctured plane.
Hence the only curves
are the two disjoint axes
in (C.HB.4), or the single
axis in (C.HB.5).
Therefore an axis not
contained in that boundary
is disjoint from it. The
preceding neighborhood
contradictions apply and
force every axis into the
boundary. For a secondary
Hopf surface, pull back
the neighborhood and its
canonical trivialization
along its finite unramified
primary covering. The
covering has the same
algebraic dimension:
meromorphic functions
descend through their
elementary symmetric
polynomials. A curve
outside the boundary
would lift to an axis
outside the pulled-back
boundary, contradicting
the already proved
primary calculation.
This proves the stated
descent and the theorem.
\end{proof}
